\section{插件、MCP 与 Skills}
%% ============================================================

Claude Code 的扩展能力来自三个层面：Plugins、MCP Servers 和 Skills。

\subsection{Skills：按需加载的知识（Tip \#8）}

Skills 是 Markdown 文件，按需扩展 Claude 的知识。与每次会话都加载的 CLAUDE.md 不同，Skills 只在与当前任务相关时才加载。

在 \texttt{.claude/skills/} 目录创建，或通过 \texttt{/plugin} 安装自带 Skills 的插件。适合存放：API 约定、部署流程、编码模式等偶尔需要的专业领域知识。

\subsection{选择合适的 MCP Server（Tip \#36）}

推荐的 MCP Server 起步组合：
\begin{itemize}
  \item \textbf{Playwright}：浏览器测试和 UI 验证
  \item \textbf{PostgreSQL/MySQL}：直接查询数据库 schema
  \item \textbf{Slack}：读取 bug 报告和对话上下文
  \item \textbf{Figma}：设计稿转代码
\end{itemize}

Claude Code 支持动态工具加载（dynamic tool loading），server 定义只在 Claude 需要时才加载。

\subsection{设置输出风格（Tip \#37）}

运行 \texttt{/config} 选择输出风格：
\begin{itemize}
  \item \textbf{Explanatory}：详细、步骤式
  \item \textbf{Concise}：简洁、行动导向
  \item \textbf{Technical}：精确、专业术语
\end{itemize}

也可以在 \texttt{\textasciitilde/.claude/output-styles/} 下创建自定义输出风格文件。

\subsection{本章小结}

扩展能力选择策略：频繁需要的指令放 CLAUDE.md，偶尔需要的领域知识用 Skills，外部服务集成用 MCP Server，能用 CLI 的场景优先 CLI。输出风格按团队偏好设置。

%% ============================================================
